\documentclass[12pt,a4paper]{article}

\usepackage[utf8]{inputenc}
\usepackage[T1]{fontenc}
\usepackage[brazilian]{babel}
\usepackage{geometry}
\usepackage{booktabs}
\usepackage{multirow}
\usepackage{array}
\usepackage{amsmath}
\usepackage{longtable}
\usepackage{graphicx}
\usepackage{float}
\usepackage{hyperref}
\usepackage{microtype}

\geometry{left=3cm,right=2cm,top=3cm,bottom=2cm}
\setlength{\parindent}{1.25cm}
\setlength{\parskip}{0.15cm}
\renewcommand{\contentsname}{Sumário}
\renewcommand{\listtablename}{Lista de Tabelas}
\renewcommand{\tablename}{Tabela}
\renewcommand{\figurename}{Figura}

\begin{document}
\begin{titlepage}
\centering
{\large Universidade Federal de Ouro Preto -- UFOP\par}
{\large Instituto de Ciências Exatas e Biológicas -- ICEB\par}
{\large Departamento de Computação -- DECOM\par}
{\large Ciência da Computação\par}
\vspace{3cm}
{\LARGE\bfseries Trabalho Prático 1\par}
\vspace{0.4cm}
{\Large BCC203 -- Estruturas de Dados II\par}
\vspace{2cm}
{\large
Stephany Fernandes Torres (25.1.4022)\par
Larissa Stefany Moura Silva (25.1.4029)\par
Lucca Sales de Souza Teodoro (21.1.4011)\par
Rafael Ferreira Lopes (25.1.4025)\par
Maria Luiza de Almeida Guerra (25.1.4024)\par
Matheus Duarte Souza Gonçalves (25.2.4152)\par
\vspace{0.6cm}
Professor: Guilherme Tavares de Assis\par}
\vfill
{\large Ouro Preto\par}
{\large 15 de Outubro de 2026\par}
\end{titlepage}

\tableofcontents
\listoftables
\newpage

\section{Introdução: Definição do Problema}

Este trabalho prático tem como objetivo implementar na linguagem C e analisar o desempenho de quatro métodos de pesquisa em memória externa: acesso sequencial indexado, árvore binária de pesquisa adequada à memória externa, árvore B e árvore B*.

A análise experimental avaliará a eficiência de cada algoritmo na manipulação de arquivos binários contendo grandes volumes de dados. O desempenho será medido considerando três pontos fundamentais: o número de transferências (leituras) da memória externa para a interna, o número de comparações entre as chaves de pesquisa e o tempo de execução.

Os resultados obtidos permitirão comparar o comportamento de cada método na prática, evidenciando suas vantagens, limitações e viabilidade em diferentes cenários de quantidade de dados no disco.

\section{Estrutura dos Dados e Metodologia de Testes}

Os testes serão feitos utilizando arquivos binários com volumes de dados: 100, 200, 2.000, 20.000, 200.000 e 2.000.000 de registros. Cada registro do arquivo terá um tamanho fixo e será composto por: uma chave de pesquisa (valor inteiro), dado1 (valor inteiro longo), dado2 (cadeia de 5000 caracteres). Tem o gerador de arquivos que tem a finalidade de gerar os arquivos para os testes. 

Esses arquivos vão ser organizados de três formas diferentes: ordenados de forma ascendente, ordenados de forma descendente e desordenados aleatoriamente. Vale ressaltar que, de acordo com as restrições de arquitetura de cada método de pesquisa implementado, algumas situações de ordenação do arquivo não se aplicam e serão desconsideradas nos testes.

Para cada cenário de teste (combinação de método, quantidade de registros e situação do arquivo), ocorrerá a pesquisa automática de 20 chaves distintas, previamente selecionadas e garantidamente presentes no arquivo. Esta abordagem visa obter uma média consistente e representativa de cada um dos quesitos avaliados (transferências, comparações e tempo).

\section{Implementação dos Métodos de Pesquisa Externa}

\subsection{Acesso Sequencial Indexado}

O método de Acesso Sequencial Indexado (ASI) implementado exige que o arquivo esteja
ordenado de forma ascendente pela chave de pesquisa. Os registros permanecem no arquivo
binário, agrupados logicamente em páginas de quatro itens, enquanto um índice esparso é
mantido na memória principal. Para cada página, o índice armazena a chave de seu primeiro
registro e a posição lógica em que a página começa no arquivo.

A pesquisa ocorre em duas etapas. Inicialmente, a tabela de índices é percorrida
sequencialmente até que seja encontrada a primeira chave maior que a chave procurada. A
página candidata é, então, a página indicada pela entrada imediatamente anterior. Em
seguida, somente essa página é lida do arquivo e seus registros são examinados
sequencialmente. Portanto, em uma pesquisa bem-sucedida ocorre apenas uma transferência
de página da memória externa para a memória interna. Em contrapartida, como a busca no
índice também é sequencial, o número de comparações cresce linearmente com a quantidade
de páginas.

\subsubsection{Funções Implementadas}

\texttt{gerarArquivo()}: cria o arquivo binário \texttt{registros.bin} antes da
execução do método. Na situação 1, aplicável ao ASI, são geradas as chaves de 1 até
$N$ em ordem crescente. O campo \texttt{dado1} recebe um valor pseudoaleatório e o
campo \texttt{dado2} é preenchido por uma cadeia de caracteres. O tempo de geração
desse arquivo não faz parte do tempo do método.

\texttt{criaTabelaIndices()}: determina a quantidade de registros a partir do tamanho
do arquivo e calcula o número de páginas por meio de
$\lceil N/\texttt{ITENS\_POR\_PAGINA}\rceil$. Em seguida, aloca dinamicamente uma
entrada do tipo \texttt{Indice} para cada página. A função lê somente o primeiro
registro de cada página, guarda sua chave e a posição lógica do início da página, e
salta os outros registros com \texttt{fseek()}. Cada primeiro registro lido é contado
como uma transferência. A última página, caso incompleta, também recebe uma entrada no
índice.

\texttt{pesquisaSequencialIndexada()}: percorre linearmente o índice, contabilizando
uma comparação para cada chave examinada, e seleciona a última página cuja chave
inicial seja menor ou igual à chave procurada. A posição lógica é convertida em um
deslocamento em bytes, e a página candidata é carregada com uma única chamada a
\texttt{fread()}. Por fim, a função percorre sequencialmente os registros lidos,
contabiliza as comparações e, se encontrar a chave, copia o registro completo para a
variável de resultado.

\texttt{executarMetodo1()}: coordena a execução do ASI. A função rejeita as situações
2 e 3, pois o método depende da ordenação ascendente; abre o arquivo; solicita a
criação da tabela; executa a pesquisa; apresenta o registro encontrado; e libera a
memória alocada. Quando a opção \texttt{-P} é utilizada, também percorre o arquivo para
mostrar suas chaves.

\texttt{sortearChaves()}: no modo experimental, acionado quando a chave informada é
$-1$, seleciona no arquivo 20 chaves distintas e existentes. As métricas das pesquisas
são acumuladas para o cálculo das médias apresentadas pelo programa principal.

\paragraph{Parâmetros de configuração}
\begin{itemize}
    \item \texttt{ITENS\_POR\_PAGINA = 4}: quantidade de registros carregados em uma página.
    \item \texttt{Indice.chave}: primeira chave da página.
    \item \texttt{Indice.primeiroRegistro}: posição lógica do primeiro registro da página no arquivo.
\end{itemize}

Não existe um limite fixo como \texttt{MAXTABELA}: o vetor de índices é alocado
dinamicamente de acordo com o tamanho do arquivo. Para o maior experimento, com
2.000.000 de registros, são necessárias
\[
\left\lceil\frac{2.000.000}{4}\right\rceil=500.000
\text{ entradas no índice}.
\]

Na construção do índice são realizadas $\lceil N/4\rceil$ transferências e nenhuma
comparação entre chaves. Na pesquisa, uma chave existente exige uma transferência de
página. No pior caso, a busca realiza até $\lceil N/4\rceil$ comparações no índice e
até quatro comparações dentro da página. Assim, a complexidade de pesquisa da
implementação é $O(N/4+4)$, equivalente a $O(N)$ em número de comparações, embora o
número de transferências da pesquisa permaneça constante.

\newpage

\subsection{Árvore Binária de Pesquisa em Memória Externa}

Neste método, a árvore binária de pesquisa é armazenada em um arquivo binário próprio
(\texttt{abp.bin}), separado do arquivo de registros. Cada nó da árvore é gravado como
um bloco do tipo \texttt{TipoNodo}, que contém o registro completo (chave, \texttt{dado1}
e \texttt{dado2}) e dois apontadores, \texttt{menor} e \texttt{maior}, que guardam a
posição no arquivo do filho à esquerda e do filho à direita. O valor $-1$ indica que o
filho não existe. A raiz da árvore é sempre o primeiro nó do arquivo (posição 0).

Diferentemente das árvores B e B*, a árvore binária de pesquisa não possui mecanismo de
balanceamento, e o seu formato depende diretamente da ordem em que os registros são
inseridos. Em arquivos ordenados (ascendente ou descendente), cada nova chave é sempre
inserida do mesmo lado, e a árvore degenera em uma lista encadeada. Isso faz com que
o número de leituras e de comparações cresça muito, tanto na construção quanto na
pesquisa.

\subsection{Funções implementadas}

\texttt{inicializaABP()}: realiza o pré-processamento do método. A função cria o arquivo
\texttt{abp.bin} e lê o arquivo de registros em blocos de \texttt{TAM\_PAGINA} (4)
registros por vez, contabilizando uma transferência por bloco lido. Cada registro lido é
transformado em um nó, com os dois apontadores iniciados em $-1$, e inserido na árvore
por meio de \texttt{insereABP()}. Ao final, é registrado o tempo total da construção.

\texttt{insereABP()}: insere um nó na árvore armazenada em arquivo. O novo nó é gravado
primeiro no final do arquivo, e sua posição é calculada a partir do tamanho do arquivo.
Se ele for o primeiro nó, torna-se a raiz. Caso contrário, a função percorre a árvore a
partir da raiz, lendo um nó por vez do arquivo (cada leitura conta como uma
transferência) e comparando as chaves: segue para o filho \texttt{menor} se a nova chave
for menor e para o filho \texttt{maior} caso contrário. Ao encontrar um apontador igual
a $-1$, atualiza esse apontador com a posição do novo nó e regrava o nó pai no arquivo.

\texttt{buscaABP()}: realiza a pesquisa de uma chave. A busca começa na raiz e, a cada
passo, um nó é lido do arquivo para a memória principal, contabilizando uma
transferência. Se a chave do nó for igual à procurada, o registro é copiado e devolvido.
Caso contrário, a busca segue pelo filho \texttt{menor} ou \texttt{maior}, conforme a
comparação. Se o apontador a ser seguido for $-1$, a chave não existe na árvore.

\texttt{executarMetodo2()}: coordena a execução do método. Abre o arquivo de registros,
chama \texttt{inicializaABP()} para construir a árvore e \texttt{buscaABP()} para
pesquisar a chave, mantendo as métricas do pré-processamento e da pesquisa separadas.
Em seguida, exibe o registro encontrado e as métricas de cada etapa. Quando o argumento
\texttt{-P} é informado, também exibe as chaves do arquivo de registros.

\newpage

\section{Árvore B}

A árvore B é uma árvore de pesquisa balanceada em que cada nó, chamado de página,
armazena várias chaves em ordem crescente. Cada página possui até $2M$ registros e
$2M+1$ apontadores para páginas filhas, e todas as folhas ficam no mesmo nível. Por
isso, a altura da árvore é pequena, da ordem de $O(\log_M n)$, mesmo para uma grande
quantidade de registros.

Neste trabalho, a árvore B é mantida em memória principal e funciona como um índice
para o arquivo de registros. Cada entrada da árvore (\texttt{TipoRegistro}) guarda apenas
a chave e a posição do registro correspondente no arquivo (\texttt{posicaoNoArquivo}).
Os registros completos permanecem no arquivo e só são lidos quando a chave procurada é
encontrada.

Na inserção, quando uma página cheia recebe um novo registro, ela é dividida em duas:
o registro do meio sobe para a página pai, e os demais são repartidos entre as duas
páginas. Se a divisão chegar à raiz, uma nova raiz é criada e a árvore cresce um nível.
Assim, a árvore cresce sempre para cima e permanece balanceada, independentemente da
ordem em que as chaves são inseridas.

\newpage

\subsection{Funções implementadas}

\texttt{CriaArvore()}: realiza o pré-processamento do método. A função percorre o
arquivo de registros do início ao fim, lendo um registro por vez (cada leitura conta
como uma transferência). Para cada registro, monta uma entrada com a chave e a posição
do registro no arquivo e a insere na árvore por meio de \texttt{Insere()}. Ao final,
registra o tempo total da construção.

\texttt{Insere()}: insere uma entrada na árvore a partir da raiz, chamando
\texttt{Ins()}. Se a raiz for dividida durante a inserção, cria uma nova página raiz
com o registro que subiu, tendo como filhos a raiz antiga e a nova página gerada pela
divisão.

\texttt{Ins()}: função recursiva que desce pela árvore até a folha onde a chave deve
ser inserida, comparando a chave com as chaves de cada página. Na volta da recursão, se
a página filha foi dividida, o registro que subiu é inserido na página atual. Se a
página atual tiver espaço, a inserção termina; caso contrário, ela também é dividida e
o registro do meio é repassado para o nível de cima. Chaves repetidas não são inseridas.

\texttt{InsereNaPagina()}: insere um registro e seu apontador à direita dentro de uma
página, deslocando as chaves maiores uma posição para a direita, de modo que a página
continue ordenada.

\texttt{Pesquisa()}: realiza a pesquisa de uma chave. Em cada página, as chaves são
percorridas até encontrar a chave procurada ou a primeira chave maior que ela. Se a
chave for encontrada, a posição do registro no arquivo é devolvida; caso contrário, a
busca continua recursivamente na página filha correspondente. Se chegar a um apontador nulo, a chave não existe na árvore.

\texttt{executarMetodo3()}: coordena a execução do método. Chama \texttt{CriaArvore()} para construir o índice e \texttt{Pesquisa()} para localizar a chave, mantendo as métricas de cada etapa separadas. Quando a chave é encontrada, o registro completo é lido do arquivo com uma única transferência, usando a posição guardada na árvore. Em seguida, exibe o registro e as métricas. Com o argumento \texttt{-P}, também exibe as chaves do arquivo de registros.

\texttt{Imprime()}: percorre a árvore em ordem e exibe as chaves armazenadas, sendo usada para verificar se a árvore foi construída corretamente.

\newpage

\section{Árvore B*}

A árvore B* é uma variação da árvore B em que existem dois tipos de página. As páginas
internas guardam apenas chaves e apontadores e servem somente como guia para a
pesquisa. Já as páginas externas (folhas) guardam todos os registros. Assim, toda
chave existente aparece em uma folha, e toda pesquisa termina em uma folha. As chaves
das páginas internas são cópias de chaves das folhas e indicam para qual lado a
pesquisa deve seguir: se a chave procurada for menor que a chave da página interna,
segue-se para a esquerda; se for maior ou igual, para a direita.

Neste trabalho, a árvore B* é mantida em memória principal e funciona como um índice
para o arquivo de registros. Cada entrada das folhas (\texttt{RegistroIndiceBEstrela})
guarda a chave e a posição do registro correspondente no arquivo. Os registros
completos permanecem no arquivo e só são lidos quando a chave procurada é encontrada.

Além disso, foi implementada a política de irmãos nas folhas, característica da árvore
B*: quando uma folha fica cheia, antes de criar uma nova página definitivamente, o
algoritmo tenta redistribuir os registros com uma folha irmã que ainda tenha espaço.
Se as irmãs também estiverem cheias, duas folhas cheias são divididas em três. Com
isso, as folhas ficam mais preenchidas, e o número de páginas criadas diminui. Nas
páginas internas, a divisão segue o funcionamento da árvore B convencional.

\newpage

\subsection{Funções implementadas}

\texttt{CriaArvoreB\_estrela()}: realiza o pré-processamento do método. A função
percorre o arquivo de registros do início ao fim, lendo um registro por vez (cada
leitura conta como uma transferência). Para cada registro, monta uma entrada com a
chave e a posição do registro no arquivo e a insere na árvore por meio de
\texttt{InsereArvoreB\_estrela()}. Ao final, registra o tempo total da construção.

\texttt{InsereArvoreB\_estrela()}: insere uma entrada na árvore a partir da raiz. Se a
árvore estiver vazia, cria uma folha como raiz. Caso contrário, chama
\texttt{insRecEstrela()} e, se a raiz for dividida, cria uma nova raiz interna com a
chave separadora e as duas páginas resultantes como filhas.

\texttt{insRecEstrela()}: função recursiva que desce pelas páginas internas até a folha
onde a chave deve ser inserida. Se a folha tiver espaço, o registro é inserido
diretamente. Se estiver cheia, a folha é dividida em duas, e a primeira chave da folha
da direita é copiada para a página pai como chave separadora. Na volta da recursão, a
página interna tenta resolver a divisão da folha com a política de irmãos; se não for
possível, insere a chave separadora em si mesma ou, se também estiver cheia, divide-se
e promove a sua chave central para o nível de cima.

\texttt{trataIrmaFolha()}: aplica a política de irmãos após a divisão de uma folha.
Primeiro tenta redistribuir os registros com a irmã direita e, depois, com a irmã
esquerda, quando alguma delas tem espaço livre, eliminando a página criada. Se as duas
irmãs estiverem cheias e a página pai tiver espaço, reparte o conteúdo das duas folhas
cheias e do novo registro entre três folhas. Em todos os casos, as chaves separadoras
da página pai são atualizadas.

\texttt{inserePaginaExterna()} e \texttt{InserePaginaInt()}: inserem, respectivamente,
um registro em uma folha e uma chave com seu apontador em uma página interna,
deslocando os elementos maiores para manter a página ordenada. A inserção em folha
também impede chaves repetidas.

\texttt{PesquisaRecEstrela()}: percorre as páginas internas comparando a chave
procurada com as chaves separadoras até chegar à folha correspondente. Na folha, as
chaves são percorridas em ordem até encontrar a chave procurada ou uma chave maior que
ela. Se a chave for encontrada, devolve a posição do registro no arquivo.

\texttt{PesquisarChaveB\_estrela()}: realiza a pesquisa completa. Chama
\texttt{PesquisaRecEstrela()} para localizar a chave no índice e, se ela for
encontrada, lê o registro completo no arquivo com uma única transferência, usando a
posição guardada na folha.

\texttt{executarMetodo4()}: coordena a execução do método, mantendo as métricas do
pré-processamento e da pesquisa separadas. Exibe o registro encontrado e as métricas
de cada etapa e, com o argumento \texttt{-P}, também exibe as chaves do arquivo de
registros. Ao final, libera a memória usada pela árvore com
\texttt{liberaArvoreB\_estrela()}.

\texttt{printArvoreB\_estrela()}: exibe a árvore por níveis, indicando quais páginas são
internas e quais são folhas, e é usada para verificar se a árvore foi construída
corretamente.

\section{Metodologia de Testes e Coleta de Resultados}

Para a realização da análise de desempenho dos métodos, os testes foram feitos de acordo com as diretrizes da segunda fase do trabalho prático. Foram gerados arquivos contendo diferentes volumes de registros ($N \in \{100, 200, 2.000, 20.000, 200.000, 2.000.000\}$), avaliados sob as situações de ordenação ascendente, descendente e aleatória.

A avaliação do tempo gasto para a criação do arquivo de dados original não foi contabilizada no tempo de execução dos métodos.

\subsection{Automação das Buscas e Cálculo de Médias}

Para cada teste (mesmo método, quantidade de registros e situação de ordem), o sistema realizou a busca automática de 20 chaves de pesquisa distintas, escolhidas aleatoriamente, mas com garantia de existência prévia no arquivo de dados.

Para cada uma das 20 buscas efetuadas, foram coletados os seguintes quesitos:
\begin{itemize}
    \item Número de transferências (leitura) de itens da memória externa para a interna.
    \item Número de comparações entre as chaves de pesquisa.
    \item Tempo de execução da pesquisa (tempo de término subtraído do tempo de início).
\end{itemize}

Os resultados finais apresentados nas seções a seguir representam a média aritmética dos valores obtidos nessas 20 execuções para cada cenário. Essa abordagem permitiu extrair uma métrica de desempenho mais estável e representativa para o comportamento de cada estrutura de dados em memória externa.

\section{Análise de Desempenho Individual}

\subsection{Avaliação: Acesso Sequencial Indexado}

A Tabela~\ref{tab:asi} apresenta os resultados obtidos na etapa de pesquisa do método
de Acesso Sequencial Indexado. Os valores correspondem à média de 20 buscas por chaves
distintas, existentes e distribuídas pelo arquivo.

\textbf{IMPORTANTE:} O método de Acesso Sequencial Indexado pressupõe arquivos ordenados pela chave de pesquisa. Dessa forma, apenas a situação de ordenação crescente foi considerada válida para execução do método.

Como cada página contém quatro registros, a tabela de índices possui
$\lceil N/4\rceil$ entradas. A construção lê o primeiro registro de cada página e não
compara chaves; por isso, a quantidade de transferências dessa etapa é determinada
diretamente por $N$, conforme a Tabela~\ref{tab:asi-criacao}. O tempo de criação não é
apresentado porque não foi registrado separadamente na rodada de testes disponível;
esse valor não deve ser estimado nem confundido com o tempo de geração de
\texttt{registros.bin}.

\begin{table}[H]
\centering
\caption{Operações realizadas na criação do índice do ASI}
\label{tab:asi-criacao}
\resizebox{\textwidth}{!}{%
\begin{tabular}{lrrrrrr}
\toprule
Operação & 100 & 200 & 2.000 & 20.000 & 200.000 & 2.000.000\\
\midrule
Entradas do índice & 25 & 50 & 500 & 5.000 & 50.000 & 500.000\\
Transferências & 25 & 50 & 500 & 5.000 & 50.000 & 500.000\\
Comparações entre chaves & 0 & 0 & 0 & 0 & 0 & 0\\
\bottomrule
\end{tabular}}
\end{table}

\begin{table}[H]
\centering
\caption{Desempenho do método de Acesso Sequencial Indexado}
\label{tab:asi}
\resizebox{\textwidth}{!}{%
\begin{tabular}{llrrrrrl}
\toprule
Ordem & Operação & 100 registros & 200 registros & 2000 registros & 20000 registros & 200000 registros & 2000000 registros\\
\midrule
\multirow{3}{*}{Ascendente} & Pesquisa (segundos) & 0.000008& 0.000004& 0.000005& 0.000010&0.000052 &0.000819\\
& Transferências & 1 & 1 & 1 & 1 & 1 & 1\\
& Comparações & 16& 29& 285& 2286&25640 &253576\\
\midrule
\multirow{3}{*}{Descendente} & Pesquisa (segundos) & -- & -- & -- & -- & -- & --\\
& Transferências & -- & -- & -- & -- & -- & --\\
& Comparações & -- & -- & -- & -- & -- & --\\
\midrule
\multirow{3}{*}{Aleatória} & Pesquisa (segundos) & -- & -- & -- & -- & -- & --\\
& Transferências & -- & -- & -- & -- & -- & --\\
& Comparações & -- & -- & -- & -- & -- & --\\
\bottomrule
\end{tabular}}
\end{table}

Os resultados mostram que todas as pesquisas bem-sucedidas utilizaram exatamente uma
transferência, independentemente do tamanho do arquivo. Esse comportamento ocorre
porque o índice fica na memória principal e somente a página candidata é carregada do
arquivo. A quantidade de comparações, entretanto, aumentou de 16 para 253.576 quando o
arquivo passou de 100 para 2.000.000 de registros. Esse crescimento aproximadamente
linear decorre da varredura sequencial do índice, que possui uma entrada para cada
quatro registros, e não da leitura da página, limitada a quatro comparações.

O tempo médio também aumentou no maior conjunto, chegando a 0,000819 segundo. As
pequenas oscilações observadas entre 100, 200 e 2.000 registros são compatíveis com a
resolução do relógio e com variações normais de execução, pois os tempos medidos são
muito pequenos. As situações descendente e aleatória foram corretamente marcadas como
não aplicáveis, uma vez que a escolha da página pressupõe que as chaves iniciais das
páginas estejam em ordem crescente.

\subsection{Avaliação: Árvore Binária de Pesquisa}

\textbf{IMPORTANTE:} O método de Árvore Binária de Pesquisa não considera entradas grandes (maiores que 10.000 itens) ordenadas previamente, pois a árvore gerada é, basicamente, uma lista encadeada, o que torna inviável a realização de testes. 

\begin{table}[H]
\centering
\caption{Desempenho do método de Árvore Binária de Pesquisa}
\label{tab:abp}
\resizebox{\textwidth}{!}{%
\begin{tabular}{llrrrrrl}
\toprule
Ordem & Operação & 100 registros & 200 registros & 2000 registros & 20000 registros & 200000 registros & 2000000 registros\\
\midrule
\multirow{3}{*}{Ascendente} & Pesquisa (segundos) & 0.000389 & 0.000626 & 0.009805 & -- & -- & --\\
& Transferências & 49 & 95 & 1064 & -- & -- & --\\
& Comparações & 146 & 285 & 3191 & -- & -- & --\\
\midrule
\multirow{3}{*}{Descendente} & Pesquisa (segundos) & 0.000257 & 0.000619 & 0.006389 & -- & -- & --\\
& Transferências & 44 & 83 & 905 & -- & -- & --\\
& Comparações & 130 & 249 & 2714 & -- & -- & --\\
\midrule
\multirow{3}{*}{Aleatória} & Pesquisa (segundos) & 0.000120 & 0.000215 & 0.000196 & 0.000324 & 0.000430 & --\\
& Transferências & 6 & 9 & 12 & 21 & 23 & --\\
& Comparações & 16 & 25 & 35 & 62 & 67 & --\\
\bottomrule
\end{tabular}}
\end{table}

\subsection{Avaliação: Árvore B}

A Tabela~\ref{tab:arvoreb} apresenta os resultados obtidos para o método da Árvore B, que mantém sua estrutura sempre balanceada na memória interna.

\begin{table}[H]
\centering
\caption{Desempenho do método da Árvore B}
\label{tab:arvoreb}
\resizebox{\textwidth}{!}{%
\begin{tabular}{llrrrrrl}
\toprule
Ordem & Operação & 100 registros & 200 registros & 2000 registros & 20000 registros & 200000 registros & 2000000 registros\\
\midrule
\multirow{3}{*}{Ascendente} & Pesquisa (segundos) & 0.000002 & 0.000002 & 0.000003 & 0.000005 & 0.000005 & 0.000054\\
& Transferências & 1 & 1 & 1 & 1 & 1 & 1\\
& Comparações & 6 & 7 & 11 & 13 & 17 & 21\\
\midrule
\multirow{3}{*}{Descendente} & Pesquisa (segundos) & 0.000002 & 0.000005 & 0.000003 & 0.000004 & 0.000005 & 0.000041\\
& Transferências & 1 & 1 & 1 & 1 & 1 & 1\\
& Comparações & 5 & 7 & 10 & 14 & 18 & 21\\
\midrule
\multirow{3}{*}{Aleatória} & Pesquisa (segundos) & 0.000002 & 0.000004 & 0.000003 & 0.000004 & 0.000004 & 0.000047\\
& Transferências & 1 & 1 & 1 & 1 & 1 & 1\\
& Comparações & 6 & 7 & 11 & 14 & 19 & 22\\
\bottomrule
\end{tabular}}
\end{table}

\subsection{Avaliação: Árvore B*}

A Tabela~\ref{tab:arvorebstar} apresenta os resultados obtidos para o método de Árvore B*, destacando a eficiência de dividir a informação apenas nas páginas folha.

\begin{table}[H]
\centering
\caption{Desempenho do método da Árvore B*}
\label{tab:arvorebstar}
\resizebox{\textwidth}{!}{%
\begin{tabular}{llrrrrrl}
\toprule
Ordem & Operação & 100 registros & 200 registros & 2000 registros & 20000 registros & 200000 registros & 2000000 registros\\
\midrule
\multirow{3}{*}{Ascendente} & Pesquisa (segundos) & 0.000065 & 0.000030 & 0.000029 & 0.000058 & 0.000060 & 0.000439\\
& Transferências & 1 & 1 & 1 & 1 & 1 & 1\\
& Comparações & 8 & 9 & 12 & 16 & 20 & 23\\
\midrule
\multirow{3}{*}{Descendente} & Pesquisa (segundos) & 0.000016 & 0.000046 & 0.000040 & 0.000053 & 0.000062 & 0.000254\\
& Transferências & 1 & 1 & 1 & 1 & 1 & 1\\
& Comparações & 8 & 9 & 12 & 15 & 20 & 23\\
\midrule
\multirow{3}{*}{Aleatória} & Pesquisa (segundos) & 0.000035 & 0.000031 & 0.000033 & 0.000046 & 0.000052 & 0.000139\\
& Transferências & 1 & 1 & 1 & 1 & 1 & 1\\
& Comparações & 8 & 10 & 14 & 17 & 21 & 25\\
\bottomrule
\end{tabular}}
\end{table}

\section{Comparação entre os Métodos}

\subsection{Pesquisa de 100 registros}

\begin{table}[H]
\centering
\caption{Comparação dos métodos na pesquisa de 100 registros}
\resizebox{\textwidth}{!}{%
\begin{tabular}{llrrrr}
\toprule
Ordem & Operação & Acesso Sequencial Indexado & Árvore Binária & Árvore B & Árvore B*\\
\midrule
\multirow{3}{*}{Ascendente} & Pesquisa (segundos) & 0.000008 & 0.000389 & 0.000002 & 0.000065\\
& Transferências & 1 & 49 & 1 & 1\\
& Comparações & 16 & 146 & 6 & 8\\
\midrule
\multirow{3}{*}{Descendente} & Pesquisa (segundos) & -- & 0.000257 & 0.000002 & 0.000016\\
& Transferências & -- & 44 & 1 & 1\\
& Comparações & -- & 130 & 5 & 8\\
\midrule
\multirow{3}{*}{Aleatória} & Pesquisa (segundos) & -- & 0.000120 & 0.000002 & 0.000035\\
& Transferências & -- & 6 & 1 & 1\\
& Comparações & -- & 16 & 6 & 8\\
\bottomrule
\end{tabular}}
\end{table}

\subsection{Pesquisa de 200 registros}

\begin{table}[H]
\centering
\caption{Comparação dos métodos na pesquisa de 200 registros}
\resizebox{\textwidth}{!}{%
\begin{tabular}{llrrrr}
\toprule
Ordem & Operação & Acesso Sequencial Indexado & Árvore Binária & Árvore B & Árvore B*\\
\midrule
\multirow{3}{*}{Ascendente} & Pesquisa (segundos) & 0.000004 & 0.000626 & 0.000002 & 0.000030\\
& Transferências & 1 & 95 & 1 & 1\\
& Comparações & 29 & 285 & 7 & 9\\
\midrule
\multirow{3}{*}{Descendente} & Pesquisa (segundos) & -- & 0.000619 & 0.000005 & 0.000046\\
& Transferências & -- & 83 & 1 & 1\\
& Comparações & -- & 249 & 7 & 9\\
\midrule
\multirow{3}{*}{Aleatória} & Pesquisa (segundos) & -- & 0.000215 & 0.000004 & 0.000031\\
& Transferências & -- & 9 & 1 & 1\\
& Comparações & -- & 25 & 7 & 10\\
\bottomrule
\end{tabular}}
\end{table}

\subsection{Pesquisa de 2.000 registros}

\begin{table}[H]
\centering
\caption{Comparação dos métodos na pesquisa de 2.000 registros}
\resizebox{\textwidth}{!}{%
\begin{tabular}{llrrrr}
\toprule
Ordem & Operação & Acesso Sequencial Indexado & Árvore Binária & Árvore B & Árvore B*\\
\midrule
\multirow{3}{*}{Ascendente} & Pesquisa (segundos) & 0.000005 & 0.009805 & 0.000003 & 0.000029\\
& Transferências & 1 & 1064 & 1 & 1\\
& Comparações & 285 & 3191 & 11 & 12\\
\midrule
\multirow{3}{*}{Descendente} & Pesquisa (segundos) & -- & 0.006389 & 0.000003 & 0.000040\\
& Transferências & -- & 905 & 1 & 1\\
& Comparações & -- & 2714 & 10 & 12\\
\midrule
\multirow{3}{*}{Aleatória} & Pesquisa (segundos) & -- & 0.000196 & 0.000003 & 0.000033\\
& Transferências & -- & 12 & 1 & 1\\
& Comparações & -- & 35 & 11 & 14\\
\bottomrule
\end{tabular}}
\end{table}

\subsection{Pesquisa de 20.000 registros}

\begin{table}[H]
\centering
\caption{Comparação dos métodos na pesquisa de 20.000 registros}
\resizebox{\textwidth}{!}{%
\begin{tabular}{llrrrr}
\toprule
Ordem & Operação & Acesso Sequencial Indexado & Árvore Binária & Árvore B & Árvore B*\\
\midrule
\multirow{3}{*}{Ascendente} & Pesquisa (segundos) & 0.000010 & -- & 0.000005 & 0.000058\\
& Transferências & 1 & -- & 1 & 1\\
& Comparações & 2286 & -- & 13 & 16\\
\midrule
\multirow{3}{*}{Descendente} & Pesquisa (segundos) & -- & -- & 0.000004 & 0.000053\\
& Transferências & -- & -- & 1 & 1\\
& Comparações & -- & -- & 14 & 15\\
\midrule
\multirow{3}{*}{Aleatória} & Pesquisa (segundos) & -- & 0.000324 & 0.000004 & 0.000046\\
& Transferências & -- & 21 & 1 & 1\\
& Comparações & -- & 62 & 14 & 17\\
\bottomrule
\end{tabular}}
\end{table}

\subsection{Pesquisa de 200.000 registros}

\begin{table}[H]
\centering
\caption{Comparação dos métodos na pesquisa de 200.000 registros}
\resizebox{\textwidth}{!}{%
\begin{tabular}{llrrrr}
\toprule
Ordem & Operação & Acesso Sequencial Indexado & Árvore Binária & Árvore B & Árvore B*\\
\midrule
\multirow{3}{*}{Ascendente} & Pesquisa (segundos) & 0.000052 & -- & 0.000005 & 0.000060\\
& Transferências & 1 & -- & 1 & 1\\
& Comparações & 25640 & -- & 17 & 20\\
\midrule
\multirow{3}{*}{Descendente} & Pesquisa (segundos) & -- & -- & 0.000005 & 0.000062\\
& Transferências & -- & -- & 1 & 1\\
& Comparações & -- & -- & 18 & 20\\
\midrule
\multirow{3}{*}{Aleatória} & Pesquisa (segundos) & -- & 0.000430 & 0.000004 & 0.000052\\
& Transferências & -- & 23 & 1 & 1\\
& Comparações & -- & 67 & 19 & 21\\
\bottomrule
\end{tabular}}
\end{table}

\subsection{Pesquisa de 2.000.000 de registros}

\begin{table}[H]
\centering
\caption{Comparação dos métodos na pesquisa de 2.000.000 de registros}
\resizebox{\textwidth}{!}{%
\begin{tabular}{llrrrr}
\toprule
Ordem & Operação & Acesso Sequencial Indexado & Árvore Binária & Árvore B & Árvore B*\\
\midrule
\multirow{3}{*}{Ascendente} & Pesquisa (segundos) & 0.000819 & -- & 0.000054 & 0.000439\\
& Transferências & 1 & -- & 1 & 1\\
& Comparações & 253576 & -- & 21 & 23\\
\midrule
\multirow{3}{*}{Descendente} & Pesquisa (segundos) & -- & -- & 0.000041 & 0.000254\\
& Transferências & -- & -- & 1 & 1\\
& Comparações & -- & -- & 21 & 23\\
\midrule
\multirow{3}{*}{Aleatória} & Pesquisa (segundos) & -- & -- & 0.000047 & 0.000139\\
& Transferências & -- & -- & 1 & 1\\
& Comparações & -- & -- & 22 & 25\\
\bottomrule
\end{tabular}}
\end{table}

\section{Dificuldades de Implementação}

Durante o desenvolvimento do trabalho, algumas dificuldades foram encontradas na
implementação e na execução dos métodos de pesquisa externa. Uma das principais
limitações esteve relacionada ao grande volume de dados utilizado nos testes,
especialmente com os arquivos de até 2.000.000 de registros. Como cada registro possui
um campo de 5000 caracteres, esses arquivos chegam a ocupar cerca de 10~GB, o que
aumentou significativamente o tempo de processamento e o espaço em disco necessário.
Além disso, como as árvores B e B* mantêm o índice em memória principal, o consumo de
memória cresce junto com a quantidade de registros, e alguns computadores utilizados
pelos integrantes do grupo apresentaram limitações na hora de rodar os testes.

No Acesso Sequencial Indexado, foi necessário tratar corretamente a última página
quando a quantidade de registros não é múltipla de quatro, além de calcular os
deslocamentos de arquivo em bytes a partir da posição lógica armazenada no índice. A
alocação dinâmica da tabela eliminou a dependência de um limite fixo, mas fez com que
seu consumo de memória crescesse proporcionalmente ao número de páginas. Outro cuidado
foi distinguir as métricas da construção do índice das métricas da pesquisa: a criação
lê a primeira chave de cada página, enquanto cada busca lê somente a página candidata.
Os testes também evidenciaram uma limitação da escolha de percorrer o índice
sequencialmente, pois ela mantém baixo o número de transferências, mas faz o número de
comparações crescer linearmente.

Outra dificuldade relevante foi a implementação da Árvore B*, considerada o método mais
complexo do trabalho. A separação entre páginas internas e folhas, a redistribuição de
registros entre folhas irmãs e a divisão de duas folhas em três exigiram bastante
cuidado, principalmente para manter corretas as chaves separadoras das páginas
internas após cada operação.

Na Árvore B, o código apresentado em aula serviu como base, mas precisou ser corrigido
e adaptado, como nas funções de inicialização e de inserção, para funcionar com os
registros do trabalho e guardar a posição de cada registro no arquivo. Já na Árvore
Binária de Pesquisa, como cada nó é lido e regravado diretamente no arquivo, foi
preciso controlar com atenção as posições de leitura e escrita, e a degeneração da
árvore com arquivos ordenados tornou a sua construção muito lenta para as maiores
quantidades de registros.

Por fim, como cada integrante ficou responsável por um método, foi necessário combinar
um mesmo critério para contar as transferências e as comparações, de forma que os
resultados pudessem ser comparados de maneira justa, além de organizar a integração do
código de todos no mesmo repositório com o Git.

\section{Conclusão}

Colocar a mão na massa e analisar os resultados deste trabalho nos permitiu entender,
na prática, onde cada método de pesquisa em memória externa brilha e onde deixa a
desejar. Para que essa comparação fosse justa, cada experimento foi feito com a
pesquisa de 20 chaves distintas, bem diferenciadas e existentes no arquivo, e os
resultados apresentados são a média dessas 20 pesquisas. Isso é importante porque o
custo de uma pesquisa depende muito de onde a chave está: uma chave perto da raiz é
encontrada rapidinho, enquanto uma chave lá no fundo da estrutura exige mais leituras e
comparações. Com a média de 20 chaves espalhadas pelo arquivo, os números mostram o
comportamento geral de cada método, e não um caso de sorte ou de azar.

O Acesso Sequencial Indexado (ASI) se destacou pelo número constante de acessos ao
disco durante a pesquisa: como a tabela de índices permanece na memória principal,
somente a página candidata precisa ser carregada. Nos experimentos, todas as buscas
bem-sucedidas realizaram uma transferência. Contudo, a implementação percorre a tabela
de índices sequencialmente; por isso, o número médio de comparações aumentou de 16,
com 100 registros, para 253.576, com 2.000.000 de registros. Uma busca binária no
índice poderia reduzir esse custo de $O(N)$ para $O(\log N)$ sem aumentar a quantidade
de páginas lidas. Além disso, o método exige um arquivo ordenado ascendentemente e é
pouco adequado a conjuntos muito dinâmicos, nos quais inserções e remoções frequentes
podem exigir reorganização do arquivo e reconstrução do índice.

Já a Árvore Binária de Pesquisa confirmou o que a teoria já avisava: ela não foi feita
para lidar com memória secundária. Como cada nó guarda um único registro e é lido
direto do arquivo, toda visita a um nó custa uma transferência, e sem nenhum mecanismo
de balanceamento o número de leituras e comparações vai às alturas. Para piorar, quando
o arquivo está ordenado, a árvore se degenera e vira basicamente uma lista encadeada.
Isso até nos impediu de rodar os testes com arquivos ordenados acima de
\textit{10.000} registros, deixando claro que não dá para contar com ela para grandes
volumes de dados.

Em contrapartida, a Árvore B e a Árvore B* entregaram tudo o que prometeram. Por serem
sempre balanceadas, elas lidaram bem com arquivos em todas as ordens (crescente,
decrescente e aleatória), e o número de comparações na pesquisa cresceu bem devagar,
de forma logarítmica, mesmo com o aumento da quantidade de registros. O segredo é o
agrupamento de várias chaves em cada página, o que mantém a árvore bem ``baixa''. Como
neste trabalho as duas ficam na memória principal funcionando como índice, encontrar
uma chave custou apenas uma leitura no disco, para buscar o registro completo no
arquivo. O preço a pagar fica no pré-processamento, já que todos os registros precisam
ser lidos e inseridos na árvore antes da primeira pesquisa.

Para fechar, vale reconhecer que a Árvore B* dá bem mais trabalho para ser
implementada. Lidar com a redistribuição de registros entre folhas irmãs e com a
divisão de duas folhas em três exige bastante cuidado. Mesmo assim, o esforço compensa,
porque as folhas ficam mais cheias e menos páginas são criadas. No fim das contas,
tanto a Árvore B quanto a B* tiveram o desempenho mais estável entre os métodos,
mostrando que apostar em estruturas multicaminho balanceadas é o caminho certo para
manipular grandes volumes de dados com eficiência.
\end{document}
